\documentclass[11pt]{article}
\usepackage[a4paper, top=18mm, bottom=18mm, left=20mm, right=20mm]{geometry}
\usepackage[T2A]{fontenc}
\usepackage[utf8]{inputenc}
\usepackage[ukrainian]{babel}
\usepackage{graphicx}
\usepackage[export]{adjustbox}
\usepackage{amsmath}
\usepackage{amssymb}
\graphicspath{{/app/Quiz/Scripts/2023/ProgAll/15BinTree/}{/app/Quiz/Scripts/2021/Mod2021_3/BinTree/}{/app/Quiz/Scripts/2020/Mod2020_4/BinTree/}}
\setlength{\parindent}{0pt}
\setlength{\parskip}{6pt}
\emergencystretch=3em
\pagestyle{empty}
\begin{document}
%begin
{\centering\Large\bfseries Контрольна робота\par}
\medskip
\noindent\rule{\linewidth}{0.5pt}
\smallskip
\noindent\textbf{Варіант \No\,2}\hfill \textit{ПІБ:}\,\underline{\hspace{60mm}}
\vspace{4mm}

%beginitem
1. %goodpath:Scripts/2018/Mod2018_1/01-good.txt
\textbf{Відмітити все, що є однією правильною лексемою:}
( 1 ) x_1

( 2 ) +

( 3 ) a1.3

( 4 ) x*y
\par\vspace{5mm}
%enditem

%beginitem
2. \textbf{Що буде виведено за виконання фрагменту коду?}

%code
#include <iostream>
using namespace std;

int g(int a, int b){
    int c = 78;
    if (b)
        c = 7;
    else if (a == 0)
         return 8;
    else 
        c = 5;
    return c;
}

int main(){
    cout << g(-4, -4);
    return 0;
}
%code
\par\vspace{5mm}
%enditem

%beginitem
3. \textbf{Що буде виведено за виконання фрагменту коду?}

%code
8//8/8*8
%code
\par\vspace{5mm}
%enditem

%beginitem
4. \textbf{Що буде виведено за виконання фрагменту коду?}

%code
print('R', end=' ')
h=['0','1','2','3','4']
h[:-2]='13'
print(h)
%code

\par\vspace{5mm}
%enditem

%beginitem
5. \textbf{Що буде виведено за виконання фрагменту коду?}
try:
    for e in range(11, 11, -2):
        if e<=7:
            continue
        print(e, end=';');e = -5
        if e<=7:
            break
    else:
        print(e,end=';')
    print(e)
except:
    print('Z')
\par\vspace{5mm}
%enditem

%beginitem
6. \textbf{Що буде виведено за виконання фрагменту коду?}

%code
cout << (!9==2 or !false< 9.0 or 2.0!=7);
%code
\par\vspace{5mm}
%enditem

%beginitem
7. \textbf{Що буде виведено за виконання фрагменту коду?}

%code
cout << (!9 == 2 or !false < 9.0 or 2.0 != 7);
%code
\par\vspace{5mm}
%enditem

%beginitem
8. \textbf{Що буде виведено за виконання фрагменту коду?}
%code
if (7 == 7)
    cout << "h";
else
    cout << "h";
%code

\par\vspace{5mm}
%enditem

%beginitem
9. 
Записати ВИРАЗ, значенням якого є список, що складається з цілих чисел від 15 до 2005 (включно), причому спочатку за спаданням записано всі, що діляться на 3, а потім решту за спаданням.\par\vspace{5mm}
%enditem

%beginitem
10. \textbf{Що буде виведено за виконання фрагменту коду?}

%code
#include <iostream>
using namespace std;

bool f(int n){
    cout<<"f";
    return n==1;
}

int main(){
    cout<<(f(-4) or f(0));
    return 0;
}
%code
\par\vspace{5mm}
%enditem

%beginitem
11. %goodpath:Scripts/2019/Mod2019_1/Lex/01-good.txt
\textbf{Відмітити все, що є однією правильною лексемою:}
( 1 ) """q""q"""

( 2 ) *

( 3 ) a1.3

( 4 ) [1]
\par\vspace{5mm}
%enditem

%beginitem
12. \textbf{Що буде виведено за виконання фрагменту коду?}

%code
try:
    m = 10
    while m>=7:
        m += -2
    print(m, end=';')
except:
    print('Z')

%code
\par\vspace{5mm}
%enditem

%beginitem
13. \textbf{Що буде виведено за виконання фрагменту коду?}

print('R', end=' ')
h=['0','1','2','3','4']
h.extend([1,2])
print(h)

\par\vspace{5mm}
%enditem

%beginitem
14. \textbf{Що буде виведено за виконання фрагменту коду?}

%code
cout << (9 / 9 / 9 / 9);
%code
\par\vspace{5mm}
%enditem

%beginitem
15. \textbf{Що буде виведено за виконання фрагменту коду?}

%code
print('R', end=' ')
h = ('a','b','c','d','e',0,1,2)
print(h[-8])
%code

\par\vspace{5mm}
%enditem

%beginitem
16. \textbf{Що буде виведено за виконання фрагменту коду?}

%code
#include <iostream>

int g(int a, int b){
    int c = 78;
    if (b)
        c = 7;
    else if (a == 0)
         return 8;
    else 
        c = 5;
    return c;
}

int main(){
    std::cout << g(-4, -4);
    return 0;
}
%code
\par\vspace{5mm}
%enditem

%beginitem
17. \textbf{Що буде виведено за виконання фрагменту коду?}

%code
a,b,c=7,8,5
def g(b):
    a-=2
    b=3
    c=1
    return a+b+c

a,b,c=3,6,2
print(g(b),a,b,c)
%code
\par\vspace{5mm}
%enditem

%beginitem
18. \textbf{Що буде виведено за виконання фрагменту коду?}
h=['0','1','2','3','4']; h[:-2]='13'; print(h)\par\vspace{5mm}
%enditem

%beginitem
19. \textbf{Що буде виведено за виконання фрагменту коду?}

print('R', end=' ')
m=10
while m>=7:
    m+=-2
print(m)

\par\vspace{5mm}
%enditem

%beginitem
20. 
Вкажіть значення та тип виразу:
not9==2 or notFalse<9.0 or 2.0!=7
\par\vspace{5mm}
%enditem

%beginitem
21. \textbf{Що буде виведено за виконання фрагменту коду?}

%code
"6.0j"==6
%code
\par\vspace{5mm}
%enditem

%beginitem
22. \textbf{Що буде виведено за виконання фрагменту коду?}

%code
h = ('a','b','c','d','e',0,1,2)
print(h[-8])
%code

\par\vspace{5mm}
%enditem

%beginitem
23. \textbf{Що буде виведено за виконання фрагменту коду?}

%code
m=7
while m<=10:
    m+=2
print(m)
%code

\par\vspace{5mm}
%enditem

%beginitem
24. \textbf{Що буде виведено за виконання фрагменту коду?}
def f(n):
    print('f', end='')
    return n==1

try:
    print(f(-4) or f(0))
except:
    print('ERR')
\par\vspace{5mm}
%enditem

%beginitem
25. 
$a$ є цілочисловим масивом типу $numpy.ndarray$ розмірів $3{\times}3$. Сума елементів масиву $a$ дорівнює 46. Чому дорівнює сума елементів масиву $a$ після виконання 
\begin{verbatim}
a -= 46
\end{verbatim}


Якщо код некоректний, то в якості відповіді зазначити ERROR.


\par\vspace{5mm}
%enditem

%beginitem
26. 
Значенням змінної \kw{a} є цілочисловий масив типу $numpy.ndarray$ розмірів $4{\times}4$. На скільки збільшиться сума елементів масиву \kw{a} після виконання
\begin{verbatim}
a.T[1] += 2
\end{verbatim}


Якщо код некоректний, то в якості відповіді зазначити ERROR.


\par\vspace{5mm}
%enditem

%beginitem
27. \textbf{Що буде виведено за виконання фрагменту коду?}
try:
    m = 7
    while m<=10:
        m += 2
    print(m, end=';')
except:
    print('Z')
\par\vspace{5mm}
%enditem

%beginitem
28. 
Написати програму, що запитує у користувача значення дійсних параметрів $x$ та $y$, обчислює для них значення виразу та повiдомляє введенi данi та результат обчислення.
$$\frac{\log_{10} x - 96}{(\arccos x - 0.6) (y - 99)}$$ 

\par\vspace{5mm}
%enditem

%beginitem
29. \textbf{Що буде виведено за виконання фрагменту коду?}
%code

x,y,z=4,0,5
y,x,z,x=y,7,z,x
print(x,y,z)
%code
\par\vspace{5mm}
%enditem

%beginitem
30. \textbf{Що буде виведено за виконання фрагменту коду?}

%code
m=10
while m>=7:
    m+=-2
print(m, end=' ')
%code

\par\vspace{5mm}
%enditem

%beginitem
31. \textbf{Що буде виведено за виконання фрагменту коду?}

%code
print('R', end=' ')
h = (0,1,2,3,4,5,6,7,8,9)
print(h[11::-2])
%code

\par\vspace{5mm}
%enditem

%beginitem
32. \textbf{Що буде виведено за виконання фрагменту коду?}

%code
not9==2 or notFalse<9.0 or 2.0!=7
%code
\par\vspace{5mm}
%enditem

%beginitem
33. \textbf{Що буде виведено за виконання фрагменту коду?}

%code
print(8//8/8*8)
%code
\par\vspace{5mm}
%enditem

%beginitem
34. %goodpath:Scripts/2019/Mod2019_1/Lex/03-good.txt
\textbf{Відмітити всі коректні оператори Python:}
( 1 ) print(3,2, end="2")

( 2 ) x=2.3; x=int(x)

( 3 ) return +

( 4 ) (2+2j)//2
\par\vspace{5mm}
%enditem

%beginitem
35. \textbf{Що буде виведено за виконання фрагменту коду?}
try:
    m = 7
    while m<=10:
        m += 2
    print(m, end=';')
except:
    print('ERR')
\par\vspace{5mm}
%enditem

%beginitem
36. \textbf{Що буде виведено за виконання фрагменту коду?}

%code
#include <iostream>

int g(int b){
    int y = 78;
    if (b) 
        y = 7;
    else if (b == 0)
         return 8;
    else
         y = 5;
    return y;
}

int main(){
    std::cout << g(-4);
    return 0;
}
%code

\par\vspace{5mm}
%enditem

%beginitem
37. 
Вкажіть значення та тип виразу:
4**-0.5*True
\par\vspace{5mm}
%enditem

%beginitem
38. \textbf{Що буде виведено за виконання фрагменту коду?}
%code

def f(n):
    print("f", end="");
    return n==1

print(f(-4) or f(0))
%code
\par\vspace{5mm}
%enditem

%beginitem
39. \textbf{Що буде виведено за виконання фрагменту коду?}

%code
print(9 != 2 > False < 9.0)
%code
\par\vspace{5mm}
%enditem

%beginitem
40. \textbf{Що буде виведено за виконання фрагменту коду?}
try:
    res = 4**-0.5*True
    if isinstance(res, bool):
        print(f'bool;{res}')
    elif isinstance(res, int):
        print(f'int;{res}')
    elif isinstance(res, float):
        print(f'float;{res:.2f}')
    else:
        print('???')
except:
    print('ERR')
\par\vspace{5mm}
%enditem

%beginitem
41. \textbf{Що буде виведено за виконання фрагменту коду?}
code
__b = 0

class B:
    __b = 1
    
    def __init__(self):
        self.__b = 2
        __b = 3
        B.__b = 4

obj = B()
try:
    print(__b, end=' ')
    print(obj.__b, end=' ')
    print(B.__b, end=' ')
    print(obj.__b, end=' ')
    print(obj._B__b, end=' ')
except:
    print('ERR')
code
\par\vspace{5mm}
%enditem

%beginitem
42. \textbf{Що буде виведено за виконання фрагменту коду?}

print('R', end=' ')
h=['0','1','2','3','4']
h[:-2]='13'
print(h)

\par\vspace{5mm}
%enditem

%beginitem
43. \textbf{Що буде виведено за виконання фрагменту коду?}

%code
print(9!=2!=False!=9.0)
%code
\par\vspace{5mm}
%enditem

%beginitem
44. \textbf{Що буде виведено за виконання фрагменту коду?}
%Оригинал был утерян. Требует отладки!
%code
__b = 0

class B:
    __b = 1
    
    def __init__(self):
        self.__b = 2
        __b = 3
        B.__b = 4

obj = B()
B.__b = 5
try:
    print(__b, end=' ')
    print(obj.__b, end=' ')
    print(B.__b, end=' ')
    print(obj.__b, end=' ')
    print(obj._B__b)
except:
    print('ERR')
%code
\par\vspace{5mm}
%enditem

%beginitem
45. \textbf{Що буде виведено за виконання фрагменту коду?}

print('R', end=' ')
h=('a','b','c','d','e',0,1,2)
print(h[-8])

\par\vspace{5mm}
%enditem

%beginitem
46. \textbf{Що буде виведено за виконання фрагменту коду?}

%code
print('R', end=' ')
m=10
while m>=7:
    m+=-2
print(m)
%code

\par\vspace{5mm}
%enditem

%beginitem
47. 
Змінні \kw{next} та \kw{prev} вузла списку позначають наступний та попередній за ним вузол відповідно. Починаючи з голови, у список записано послідовні цілі числа від~$-63$ до~$22$. Нехай змінна \kw{p} позначає вузол, що зберігає значення~$-4$.
Яке значення зберігається у вузлі \kw{p.next.next.next.next.next} ?
\par\vspace{5mm}
%enditem

%beginitem
48. \textbf{Що буде виведено за виконання фрагменту коду?}

%code
cout << 8 / 8 / 8 / 8;
%code
\par\vspace{5mm}
%enditem

%beginitem
49. \textbf{Що буде виведено за виконання фрагменту коду?}

print('R', end=' ')
m=7
while m<=10:
    m+=2
print(m)

\par\vspace{5mm}
%enditem

%beginitem
50. \textbf{Що буде виведено за виконання фрагменту коду?}

try:
    print(7,end="")
    print(8//0.0,end="")
    print(5,end="")
except ZeroDivisionError: print(3,end="")
except Exception: print(6,end="")
else: print(2,end="")
finally: print(1,end="")
\par\vspace{5mm}
%enditem

%beginitem
51. 
Записати ВИРАЗ, значенням якого є множина, що складається з цілих чисел від 15 до 2005 (включно), що діляться на 3.\par\vspace{5mm}
%enditem

%beginitem
52. \textbf{Що буде виведено за виконання фрагменту коду?}

%code
print(4**-0.5*True)
%code
\par\vspace{5mm}
%enditem

%beginitem
53. 
Змінні $next$ та $prev$ вузла списку позначають наступний та попередній за ним вузол відповідно. Починаючи з голови, у список записано послідовні цілі числа від $-63$ до $22$. Нехай змінна $p$ позначає вузол, що зберігає значення $-4$.
Яке значення зберігається у вузлі $p.next.next.next.next.next$ ?
\par\vspace{5mm}
%enditem

%beginitem
54. \textbf{Що буде виведено за виконання фрагменту коду?}

%code
print('R', end=' ')
h=('a','b','c','d','e',0,1,2)
print(h[-8])
%code

\par\vspace{5mm}
%enditem

%beginitem
55. \textbf{Що буде виведено за виконання фрагменту коду?}
try:
    res = 8//8/8*8
    if isinstance(res, bool):
        print(f'bool;{res}')
    elif isinstance(res, int):
        print(f'int;{res}')
    elif isinstance(res, float):
        print(f'float;{res:.2f}')
    else:
        print('???')
except:
    print('ERR')
\par\vspace{5mm}
%enditem

%beginitem
56. \textbf{Що буде виведено за виконання фрагменту коду?}
def f(arg, n):
    n = 7
    tmp = arg
    tmp[:-2] = 'bd'
    return len(tmp)>3

try:
    h = ['0','1','2','3','4']
    n = 8
    f(h, n)
    print(n, end=';')
    for el in h:
        print(el, end=';')
except:
    print('Z')
\par\vspace{5mm}
%enditem

%beginitem
57. \textbf{Що буде виведено за виконання фрагменту коду?}
try:
    res = not9==2 or notFalse<9.0 or 2.0!=7
    if isinstance(res, bool):
        print(f'bool;{res}')
    elif isinstance(res, int):
        print(f'int;{res}')
    elif isinstance(res, float):
        print(f'float;{res:.2f}')
    else:
        print('???')
except:
    print('Z')
\par\vspace{5mm}
%enditem

%beginitem
58. \textbf{Що буде виведено за виконання фрагменту коду?}
%code
a, b, c = 7, 6, 8
def g(a, b, c=9):
    print(a, b, c, end=" ")

g(b=4, c=0, a=5)
print(a, b, c)
%code
\par\vspace{5mm}
%enditem

%beginitem
59. \textbf{Що буде виведено за виконання фрагменту коду?}

%code
h=['0','1','2','3']
h.extend([1,2])
print(h)
%code

\par\vspace{5mm}
%enditem

%beginitem
60. \textbf{Що буде виведено за виконання фрагменту коду?}
try:
    for e in range(11, 11, -2):
        if e<=7:
            continue
        print(e, end=';');e = -5
        if e<=7:
            break
    else:
        print(e,end=';')
    print(e)
except:
    print('ERR')
\par\vspace{5mm}
%enditem

%beginitem
61. \textbf{Що буде виведено за виконання фрагменту коду?}
%code
a,b,c=7,6,8
def g(a,b,c=9):
    print(a,b,c,end=" ")

g(b=4,c=0,a=5)
print(a,b,c)
%code
\par\vspace{5mm}
%enditem

%beginitem
62. \textbf{Що буде виведено за виконання фрагменту коду?}

%code
try:
    m = 7
    while m<=10:
        m += 2
    print(m, end=';')
except:
    print('Z')

%code
\par\vspace{5mm}
%enditem

%beginitem
63. \textbf{Що буде виведено за виконання фрагменту коду?}

%code
h = (0,1,2,3,4,5,6,7,8,9)
print(h[11::-2])
%code

\par\vspace{5mm}
%enditem

%beginitem
64. 
Написати програму, що запитує у користувача значення дійсних параметрів $x$ та $y$, обчислює для них значення виразу 
$$\frac{\log_{10} x - 96}{(\arccos x - 0.6) (y - 99)}$$ 
та повiдомляє введенi данi та результат обчислення.

\par\vspace{5mm}
%enditem

%beginitem
65. \textbf{Що буде виведено за виконання фрагменту коду?}

%code
def g(a,b):
    c=78
    if b:
        c=7
    elif a==0:
         return 8
    else: 
        c=5
    return c

print(g(-4,-4))
%code
\par\vspace{5mm}
%enditem

%beginitem
66. \textbf{Що буде виведено за виконання фрагменту коду?}

%code
a,b,c=7,8,5
def g(b):
    a=2
    b=3
    c=1
    return a+b+c

a,b,c=3,6,2
print(g(b),a,b,c)
%code
\par\vspace{5mm}
%enditem

%beginitem
67. \textbf{Що буде виведено за виконання фрагменту коду?}
%code
if (7 == 7)
    cout << "h";
else
    cout << "y";
%code

\par\vspace{5mm}
%enditem

%beginitem
68. \textbf{Що буде виведено за виконання фрагменту коду?}

print('R', end=' ')
h=(0,1,2,3,4,5,6,7,8,9)
print(h[11::-2])

\par\vspace{5mm}
%enditem

%beginitem
69. \textbf{Що буде виведено за виконання фрагменту коду?}
try:
    m = 10
    while m>=7:
        m += -2
    print(m, end=';')
except:
    print('ERR')
\par\vspace{5mm}
%enditem

%beginitem
70. \textbf{Що буде виведено за виконання фрагменту коду?}

%code
print("6.0j" == 6)
%code
\par\vspace{5mm}
%enditem

%beginitem
71. \textbf{Що буде виведено за виконання фрагменту коду?}
try:
    res = "6.0j"==6j
    if isinstance(res, bool):
        print(f'bool;{res}')
    elif isinstance(res, int):
        print(f'int;{res}')
    elif isinstance(res, float):
        print(f'float;{res:.2f}')
    else:
        print('???')
except:
    print('ERR')
\par\vspace{5mm}
%enditem

%beginitem
72. \textbf{Що буде виведено за виконання фрагменту коду?}

%code
#include <iostream>
using namespace std;

int g(int b){
    int y = 78;
    if (b) 
        y = 7;
    else if (b == 0)
         return 8;
    else
         y = 5;
    return y;
}

int main(){
    cout << g(-4);
    return 0;
}
%code

\par\vspace{5mm}
%enditem

%beginitem
73. \textbf{Що буде виведено за виконання фрагменту коду?}

%code
cout << 9 / 9 / 9 / 9;
%code
\par\vspace{5mm}
%enditem

%beginitem
74. 
Вкажіть значення та тип виразу:
8//8/8*8
\par\vspace{5mm}
%enditem

%beginitem
75. \textbf{Що буде виведено за виконання фрагменту коду?}
try:
    res = not9==2 or notFalse<9.0 or 2.0!=7
    if isinstance(res, bool):
        print(f'bool;{res}')
    elif isinstance(res, int):
        print(f'int;{res}')
    elif isinstance(res, float):
        print(f'float;{res:.2f}')
    else:
        print('???')
except:
    print('ERR')
\par\vspace{5mm}
%enditem

%beginitem
76. 
Записати ВИРАЗ, значенням якого є словник, ключами якого є цілі числа від 15 до 2005 (включно), що діляться на 3, а ключам відповідають їх куби \par\vspace{5mm}
%enditem

%beginitem
77. \textbf{Що буде виведено за виконання фрагменту коду?}

%code
cout << (8 / 8 / 8 / 8);
%code
\par\vspace{5mm}
%enditem

%beginitem
78. 
Змінні \kw{next} та \kw{prev} вузла списку позначають наступний та попередній за ним вузол відповідно. Починаючи з голови, у список записано послідовні цілі числа від~$-63$ до~$22$. Нехай змінна \kw{p} позначає вузол, що зберігає значення~$-4$.
Яке число зберігається у вузлі \kw{p.next.next.next.next.next} ?
\par\vspace{5mm}
%enditem

%beginitem
79. \textbf{Що буде виведено за виконання фрагменту коду?}

%code
print('R', end=' ')
h=['0','1','2','3','4']
h.extend([1,2])
print(h)
%code

\par\vspace{5mm}
%enditem

%beginitem
80. 
Задано тип вузла списку
struct Node \{int n; Node *prev, *next;\};
У списку зберігаються послідовні цілі числа від -63 до 22. Вказівник p встановлено на вузол, в якому зберігається число -4.
Чому дорівнює значення p->next->next->next->next->next->n ?\par\vspace{5mm}
%enditem

%beginitem
81. \textbf{Що буде виведено за виконання фрагменту коду?}

%code
print('R', end=' ')
m=7
while m<=10:
    m+=2
print(m)
%code

\par\vspace{5mm}
%enditem

%beginitem
82. \textbf{Що буде виведено за виконання фрагменту коду?}
try:
    res = 4**-0.5*True
    if isinstance(res, bool):
        print(f'bool;{res}')
    elif isinstance(res, int):
        print(f'int;{res}')
    elif isinstance(res, float):
        print(f'float;{res:.2f}')
    else:
        print('???')
except:
    print('Z')
\par\vspace{5mm}
%enditem

%beginitem
83. \textbf{Що буде виведено за виконання фрагменту коду?}

%code
print('R', end=' ')
h=(0,1,2,3,4,5,6,7,8,9)
print(h[11::-2])
%code

\par\vspace{5mm}
%enditem

%beginitem
84. \textbf{Що буде виведено за виконання фрагменту коду?}

%code
4**-0.5*True
%code
\par\vspace{5mm}
%enditem

%beginitem
85. 
Кожен рядок текстового файлу містить по два дійсних числа, розділені білими символами.
Написати функцію, що для заданого текстового файлу будує новий текстовий файл, рядки якого крім зображень дйсних чисел ще містять результат обчислення на них функції $f$. У вихідному файлі зображення чисел мають розділятися пробілами. 
$$f(x, y) =\frac{\log_{10} x - 96}{(\arccos x - 0.6) (y - 99)}$$ 
Якість коду також оцінюється.

\par\vspace{5mm}
%enditem

%beginitem
86. 
Вкажіть значення та тип виразу:
"6.0j"==6
\par\vspace{5mm}
%enditem

%beginitem
87. \textbf{Що буде виведено за виконання фрагменту коду?}

x,y,z=4,0,5
y,x,z,x=y,7,z,x
print(x,y,z)
\par\vspace{5mm}
%enditem

%beginitem
88. \textbf{Що буде виведено за виконання фрагменту коду?}
def f(arg, n):
    n = 7
    tmp = arg
    tmp[:-2] = 'bd'
    return len(tmp)>3

try:
    h = ['0','1','2','3','4']
    n = 8
    f(h, n)
    print(n, end=';')
    for el in h:
        print(el, end=';')
except:
    print('ERR')
\par\vspace{5mm}
%enditem

%beginitem
89. \textbf{Що буде виведено за виконання фрагменту коду?}
try:
    res = 8//8/8*8
    if isinstance(res, bool):
        print(f'bool;{res}')
    elif isinstance(res, int):
        print(f'int;{res}')
    elif isinstance(res, float):
        print(f'float;{res:.2f}')
    else:
        print('???')
except:
    print('Z')
\par\vspace{5mm}
%enditem

%beginitem
90. %goodpath:Scripts/2018/Mod2018_1/03-good.txt
\textbf{Відмітити все, що є коректним оператором С++:}
( 1 ) int x,y;

( 2 ) bool r=('N'>'5');

( 3 ) "a"=6;

( 4 ) int f(int a, int 4);
\par\vspace{5mm}
%enditem

%beginitem
91. \textbf{Що буде виведено за виконання фрагменту коду?}

def g(a,b):
    c=78
    if b:
        c=7
    elif a==0:
         return 8
    else: 
        c=5
    return c

print(g(-4,-4))\par\vspace{5mm}
%enditem

%beginitem
92. 
Написати програму, що запитує у користувача значення дійсних параметрів $x$ та $y$, обчислює для них значення виразу 
$$\frac{\log_{10} x - 96}{(\arccos x - 0.6) (y - 99)}$$ 
та повiдомляє введенi данi та результат обчислення.
 Оцінюється правильність та якість коду.

\par\vspace{5mm}
%enditem

%beginitem
93. \textbf{Що буде виведено за виконання фрагменту коду?}
h=['0','1','2','3','4']; h.extend([1,2]); print(h)\par\vspace{5mm}
%enditem

%beginitem
94. \textbf{Що буде виведено за виконання фрагменту коду?}
%code

a,b,c=7,6,8
def g(a,b,c=9):
    print(a,b,c,end="")

g(b=4,c=0,a=5)
print(a,b,c)

%code
\par\vspace{5mm}
%enditem

%beginitem
95. 
В умовах попередньої задачі було виконано p->next=p->prev->next;

Чому дорівнює значення p->next->next->next->next->n ?\par\vspace{5mm}
%enditem

%beginitem
96. \textbf{Що буде виведено за виконання фрагменту коду?}
try:
    print(7, end=';')
    print(8//0.0, end=';')
    print(5, end=';')
except ValueError:
    print(3, end=';')
except TypeError:
    print(6, end=';')
except:
    print('ERR', end=';')
else:
    print(2, end=';')
finally:
    print(1)
\par\vspace{5mm}
%enditem

%beginitem
97. 

\begin{center}
	\adjustimage{max size={0.8\linewidth}{0.8\paperheight}}
	{../15BinTree/trees_exp/68.png}
\end{center}
Указати послідовність міток вершин дерева, зображеного на рис., що утвориться в результаті обходу дерева в оберненому порядку. Мітки записувати через пробіл.\par Алгоритм починає роботу з кореня (на рис. це верхня вершина).
%answer:68 оберненому
\par\vspace{5mm}
%enditem

%beginitem
98. 
Рядки журналу через = містять ім'я користувача (ідентифікатор), час події,тип події, код події, наприклад
\begin{verbatim}
s="username = 2022-05-05 13:26 = ERROR = 404"
\end{verbatim}
у коді 
\begin{verbatim}
res = re.fullmatch('???', s) 
s = ''
code_str = ??? 
\end{verbatim}
чим треба замінити кожен з ???, щоб значенням code\_str був код події? 


\par\vspace{5mm}
%enditem

%beginitem
99. \textbf{Що буде виведено за виконання фрагменту коду?}
try:
    m = 10
    while m>=7:
        m += -2
    print(m, end=';')
except:
    print('Z')
\par\vspace{5mm}
%enditem

%beginitem
100. \textbf{Що буде виведено за виконання фрагменту коду?}

%code
def g(b):
    y=78
    if b: 
        y=7
    elif b==0:
         return 8
    else:
         y=5
    return y

print(g(-4))
%code
\par\vspace{5mm}
%enditem

%end
\newpage
\end{document}

